\documentclass[12pt,reqno, twoside]{amsbook}
%%%%%%%%%%%%%%%%%%%%%%%%%%%%%%%%%%%%%%%%%%%%%%%%%%%%%%%%%%%%%%%%%%%%%%%%%%%%%%%%%%%%%%%%%%%%%%%%%%%%%%%%%%%%%%%%%%%%%%%%%%%%%%%%%%%%%%%%%%%%%%%%%%%%%%%%%%%%%%%%%%%%%%%%%%%%%%%%%%%%%%%%%%%%%%%%%%%%%%%%%%%%%%%%%%%%%%%%%%%%%%%%%%%%%%%%%%%%%%%%%%%%%%%%%%%%
\usepackage{eurosym}
\usepackage{amsmath}
\usepackage{amssymb}
\usepackage{amsfonts}
\usepackage[T1]{fontenc}
\usepackage[utf8]{inputenc}
\usepackage[onehalfspacing]{setspace}
\usepackage{chngcntr}
\usepackage{graphicx}
\graphicspath{{images/}}
\usepackage[a4paper, margin=2.5cm]{geometry}
\usepackage[portuguese]{babel}
\usepackage{fancyhdr}
\usepackage{titlesec}
\usepackage{enumitem}
\usepackage{etoolbox}
\usepackage{array}
\usepackage{booktabs}

\makeatletter
    \def\section{\@startsection{section}{1}%
      \z@{.5\linespacing\@plus.7\linespacing}{.25\linespacing}%
      {\normalfont\bfseries\flushleft}}
    \def\subsection{\@startsection{subsection}{2}%
      \z@{.5\linespacing\@plus.7\linespacing}{.25\linespacing}%
      {\normalfont\bfseries\flushleft}}

\makeatother
\setcounter{MaxMatrixCols}{10}

\providecommand{\U}[1]{\protect\rule{.1in}{.1in}}
\theoremstyle{plain}
\newtheorem{acknowledgement}{Agradecimento}
\newtheorem{algorithm}{Algoritmo}[chapter]
\newtheorem{axiom}{Axioma}[chapter]
\newtheorem{case}{Caso}[chapter]
\newtheorem{claim}{Claim}[chapter]
\newtheorem{conclusion}{Conclusão}[chapter]
\newtheorem{condition}{Condição}[chapter]
\newtheorem{conjecture}{Conjetura}[chapter]
\newtheorem{corollary}{Corolário}[chapter]
\newtheorem{criterion}{Critério}[chapter]
\newtheorem{definition}{Definição}[chapter]
\newtheorem{example}{Exemplo}[chapter]
\newtheorem{exercise}{Exercício}[chapter]
\newtheorem{lemma}{Lema}[chapter]
\newtheorem{notation}{Notação}[chapter]
\newtheorem{problem}{Problema}[chapter]
\newtheorem{proposition}{Proposição}[chapter]
\newtheorem{remark}{Nota}[chapter]
\newtheorem{solution}{Solução}[chapter]
\newtheorem{summary}{Sumário}[chapter]
\newtheorem{theorem}{Teorema}[chapter]
\numberwithin{equation}{chapter}

\newenvironment{dedication}
  {
   \thispagestyle{empty}
   \vspace*{\stretch{1}}
   \itshape             
   \raggedleft          
  }
  {\par 
   \vspace{\stretch{3}} 
  }
  \numberwithin{section}{chapter}
\fancyhead{}
\fancyfoot{}
\pagestyle{fancy}
\fancyfoot[LE,RO]{\thepage}

\makeatletter
\def\ps@plain{\ps@empty
 \def\@evenfoot{%
   \normalfont\scriptsize
   \rlap{\thepage}\hfil
   }%
 \def\@oddfoot{%
   \normalfont\scriptsize \hfil
   \llap{\thepage}}%
}
\makeatother
\renewcommand{\headrulewidth}{0pt}
\renewcommand{\footrulewidth}{0pt}

\begin{document}
\begin{titlepage}
\centering

% --- LOGÓTIPO DA INSTITUIÇÃO ---
% Se tiveres a imagem do logótipo no Overleaf/VS Code (ex: logo_iscte.png), descomenta a linha abaixo:
\includegraphics[width=0.45\textwidth]{logo_iscte_pt.png} \\[1.5cm]

{\Large \textbf{ISCTE -- Instituto Universitário de Lisboa}} \\[0.3cm]
{\large Escola de Tecnologias e Arquitetura} \\[0.2cm]
{\large Departamento de Ciências e Tecnologias da Informação} \\[1.5cm]

% --- TÍTULO DA CADEIRA / UNIDADE CURRICULAR ---
{\small \textbf{SEGURANÇA DE REDES E SISTEMAS DE INFORMAÇÃO}} \\[0.3cm]
{\small Relatório do Projeto -- Parte 1 (Teórica)} \\[1.5cm]

\rule{\linewidth}{0.5mm} \\[0.4cm]
% --- TÍTULO DO TRABALHO ---
{\Large \bfseries Middleware Adaptativo de Privacidade e Redução de Dados em Sistemas de Informação} \\[0.3cm]
{\large \textit{Implementação de um Modelo Inspirado em PADRES/PRIPRO para Proteção de Dados e Segurança}}
\rule{\linewidth}{0.5mm} \\[2.5cm]

% --- AUTOR E DOCENTE ---
\begin{minipage}{0.48\textwidth}
    \begin{flushleft} \large
        \textbf{Elaborado por:}\\[0.2cm]
        [O Teu Nome Completo] \\
        Número de Aluno: [O Teu Número]
    \end{flushleft}
\end{minipage}
\hfill
\begin{minipage}{0.48\textwidth}
    \begin{flushright} \large
        \textbf{Docente:}\\[0.2cm]
        Prof. [Nome do Professor] \\
        \textit{Iscte-IUL}
    \end{flushright}
\end{minipage}

\vfill

% --- DATA E LOCAL ---
{\large Lisboa, \today}

\end{titlepage}
\thispagestyle{empty}
\null
\newpage
\setcounter{page}{1}

\frontmatter
\addtocontents{toc}{\setcounter{tocdepth}{2}}
\thispagestyle{empty}

\chapter*{Resumo}

\renewcommand{\contentsname}{Índice}
\tableofcontents

\mainmatter

\hyphenation{middleware privacy-aware proxy anonymization pseudonimization}

\setcounter{page}{1} \pagenumbering{arabic}
\chapter{Introdução}
\noindent A crescente integração entre sistemas de informação tem aumentado a circulação de dados através de interfaces de programação de aplicações (APIs), em particular de serviços HTTP/REST. Embora esta partilha permita automatizar processos e disponibilizar informação a diferentes entidades, uma resposta pode conter atributos que não são necessários para todos os seus destinatários. A exposição desses atributos aumenta o risco para a privacidade e pode também implicar a transmissão de um volume de dados superior ao necessário.

\noindent Neste contexto, este trabalho propõe o desenvolvimento de um middleware/proxy de privacidade em Python, colocado entre os clientes e os serviços de backend. A camada intermédia deverá identificar o perfil associado ao pedido e adaptar a resposta antes de a devolver ao cliente. Deste modo, pretende-se aproximar o controlo de acesso da própria representação dos dados transmitidos, sem pressupor que todos os clientes autorizados necessitam da mesma informação.

\section{Contexto e Motivação}
\noindent Em ecossistemas de APIs RESTful, os mecanismos de autorização determinam frequentemente se um pedido pode aceder a um recurso, mas não adaptam necessariamente o conteúdo da resposta às necessidades concretas de quem o solicita. Assim, um consumidor autorizado a consultar um recurso pode receber dados pessoais identificáveis (PII) ou outros atributos sensíveis, mesmo quando apenas necessita de uma visão parcial ou de métricas consolidadas. A aplicação de uma política uniforme pode, por isso, favorecer a utilidade operacional à custa da minimização da informação divulgada.

\noindent A motivação do projeto é explorar uma abordagem intermédia que combine perfis de acesso com transformações do payload. Inspirado nos modelos PRIPRO/PRIADA, que orientam a gestão dinâmica de perfis de privacidade e a gestão de contexto em ambientes distribuídos \cite{cesconetto2020}, e PADRES, associado à redução de dados consciente da privacidade e à regulação de segurança \cite{pereira2022}, o middleware deverá adaptar a resposta a três perfis previstos na proposta: o Administrador recebe os dados necessários às suas funções e permissões; o Operacional/Analista recebe dados com mascaragem ou pseudonimização de atributos confidenciais; e o Público/Terceiros recebe informação agregada e reduzida. A aplicação destas transformações procura limitar a exposição de dados e, nos casos em que os campos são omitidos ou agregados, reduzir também o payload transmitido.

\noindent Esta abordagem exige considerar simultaneamente a proteção da informação e a sua utilidade. A transformação não deve remover dados necessários à finalidade legítima do pedido, e a redução do payload não constitui, por si só, uma garantia de anonimização ou de conformidade regulatória. Por esse motivo, o trabalho enquadra a proposta nos princípios do RGPD e na resiliência dos sistemas de informação abordada pela Diretiva NIS2, e prevê avaliar o comportamento do protótipo através de testes funcionais e métricas de desempenho.

\section{Definição do Problema}
\noindent O problema central consiste na dificuldade de conciliar a disponibilização de informação útil com a limitação dos dados pessoais e do volume de informação transmitidos a cada consumidor. Uma decisão estática de permitir ou negar o acesso, incluindo uma política baseada em papéis (RBAC), pode não distinguir entre destinatários com necessidades diferentes perante o mesmo recurso. Como consequência, mesmo quando o acesso é legítimo, a resposta pode conter mais informação do que a necessária.

\noindent Pretende-se, por isso, estudar como um middleware pode usar o perfil associado ao pedido para aplicar políticas de transformação coerentes antes de encaminhar a resposta. A solução deverá articular a identificação do perfil, a seleção dos campos a manter ou transformar e a entrega da resposta adaptada. A sua avaliação deverá verificar se cada perfil recebe o tipo de informação esperado e medir o custo introduzido pela camada intermédia, incluindo latência e consumo de recursos, bem como a variação do tamanho do payload. O desafio está em reduzir a exposição e a transmissão de dados sem comprometer a correção funcional das respostas.

\section{Objetivos do Trabalho}
\noindent O objetivo geral é conceber e avaliar um protótipo de middleware/proxy que adapte respostas de APIs segundo perfis de acesso, combinando proteção da privacidade e redução de dados. Para o concretizar, o projeto divide-se em duas componentes:
\begin{itemize}[leftmargin=*]
    \item \textbf{Objetivo Teórico (Parte 1):} Rever abordagens de controlo de acesso adaptativo e de proteção da privacidade, enquadrar a proposta face ao RGPD e à Diretiva NIS2 e formalizar métricas para discutir a informação exposta, com base na Entropia de Shannon, e a redução do payload, através de $\Delta D$.
    \item \textbf{Objetivo Prático (Parte 2):} Implementar um protótipo modular em Python, recorrendo a FastAPI ou Flask, que intercepte pedidos HTTP/REST e adapte as respostas para os perfis definidos. Validar funcionalmente as transformações e comparar o comportamento com e sem middleware, medindo a latência, o consumo de CPU e memória e o tamanho dos dados transmitidos.
\end{itemize}

\noindent O trabalho incide sobre um protótipo e a sua avaliação experimental; os resultados permitirão caracterizar a abordagem no cenário de teste, não constituindo por si só uma certificação de conformidade nem uma demonstração de adequação a ambientes de produção.
\chapter{Enquadramento Regulatório e Normativo}
\noindent A conceção de um middleware que intermedeia pedidos e respostas deve considerar tanto a proteção dos dados pessoais como a segurança e a resiliência dos sistemas de informação. Neste projeto, o Regulamento Geral sobre a Proteção de Dados (RGPD) é relevante quando os dados tratados permitem identificar diretamente ou indiretamente uma pessoa singular. A Diretiva NIS2, por sua vez, estabelece requisitos de gestão de riscos de cibersegurança para as entidades abrangidas pelo seu âmbito. Os dois instrumentos são complementares, mas têm objetivos e âmbitos de aplicação distintos: a existência do protótipo, por si só, não demonstra conformidade com nenhum deles.

\section{Regulamento Geral sobre a Proteção de Dados (RGPD)}
\noindent O RGPD estabelece princípios e obrigações aplicáveis ao tratamento de dados pessoais \cite{rgpd2016}. Para o middleware proposto, são particularmente relevantes:
\begin{itemize}[leftmargin=*]
    \item \textbf{Limitação das finalidades e minimização (Artigo 5.º, n.º 1, alíneas b) e c)):} Os dados devem ser tratados para finalidades determinadas e explícitas, e limitar-se ao que é necessário para essas finalidades. No protótipo, isto implica definir, por recurso e perfil, que campos são necessários, em vez de devolver automaticamente toda a resposta do serviço de backend.
    \item \textbf{Exatidão e limitação da conservação (Artigo 5.º, n.º 1, alíneas d) e e)):} Os dados devem ser exatos e conservados apenas durante o período necessário. Embora o middleware se concentre na resposta em trânsito, os registos e dados temporários que possa produzir também devem ser considerados na conceção e nos testes.
    \item \textbf{Integridade, confidencialidade e responsabilidade (Artigo 5.º, n.º 1, alínea f), e n.º 2):} O responsável pelo tratamento deve assegurar proteção adequada contra tratamento não autorizado ou ilícito e conseguir demonstrar o cumprimento dos princípios aplicáveis. A adaptação de respostas deve, portanto, ser acompanhada por controlos de acesso e por políticas verificáveis, não apenas por remoção de campos.
    \item \textbf{Proteção de dados desde a conceção e por defeito (Artigo 25.º):} As medidas técnicas e organizacionais devem ser integradas na conceção do tratamento e, por defeito, limitar a quantidade de dados, a extensão do tratamento, o período de conservação e a acessibilidade. A separação das respostas por perfil constitui uma possível medida técnica, desde que os perfis e as regras sejam definidos de forma adequada à finalidade.
    \item \textbf{Segurança do tratamento (Artigo 32.º):} As medidas devem ser adequadas ao risco e podem incluir pseudonimização e cifragem, mecanismos para garantir confidencialidade, integridade, disponibilidade e resiliência, bem como processos de teste e avaliação regulares. O artigo não prescreve uma tecnologia única: a escolha deve considerar o contexto, os riscos e o estado da técnica.
\end{itemize}

\noindent É importante distinguir as transformações previstas. A \textit{mascaragem} pode ocultar parte de um valor numa resposta, mas não garante, por si só, que a pessoa deixe de ser identificável. A \textit{pseudonimização}, nos termos do artigo 4.º, n.º 5, substitui identificadores por outros valores e exige que a informação adicional que permita atribuir os dados a uma pessoa seja mantida separadamente e protegida; os dados pseudonimizados continuam a ser dados pessoais quando a reidentificação é possível. A anonimização exige que a identificação deixe efetivamente de ser razoavelmente possível, o que depende também do contexto e de outros dados disponíveis. Assim, reduzir o tamanho do payload ou aplicar hashing não deve ser apresentado automaticamente como anonimização.

\section{Diretiva NIS2}
\noindent A Diretiva (UE) 2022/2555, conhecida como NIS2, procura estabelecer um nível elevado comum de cibersegurança na União Europeia \cite{nis22022}. A sua aplicação depende do âmbito definido na Diretiva e da legislação nacional que a transpõe; não se deve, por isso, assumir que qualquer serviço ou protótipo está automaticamente sujeito às mesmas obrigações. Para as entidades abrangidas, a Diretiva prevê responsabilidades de governação e medidas de gestão de riscos, tendo em conta a exposição a riscos, a dimensão da entidade e a probabilidade e gravidade dos incidentes.

\noindent O artigo 21.º identifica áreas de medidas de gestão de riscos, entre as quais o tratamento de incidentes, a continuidade das atividades, a segurança da cadeia de fornecimento, a segurança na aquisição, desenvolvimento e manutenção de sistemas, a avaliação da eficácia das medidas, políticas de criptografia, controlo de acessos e, quando adequado, autenticação multifator e comunicações seguras. A Diretiva não se reduz à proteção de dados pessoais: o seu foco inclui a disponibilidade e a resiliência dos sistemas e serviços.

\noindent No contexto do projeto, um middleware que recebe e reencaminha pedidos é também um componente cuja configuração e funcionamento devem ser considerados no modelo de ameaças. Uma política de filtragem pode reduzir a exposição de campos, mas não substitui medidas como autenticação robusta, validação dos tokens, proteção do tráfego, gestão segura de segredos, atualização de dependências, tratamento controlado de erros e preparação para falhas. A relevância e a proporcionalidade de cada medida devem ser avaliadas em função do cenário de utilização.

\section{Implicações para o protótipo}
\noindent O enquadramento regulatório traduz-se em requisitos de conceção e de validação, não numa declaração de conformidade:
\begin{itemize}[leftmargin=*]
    \item \textbf{Autorização antes da transformação:} o perfil deve ser obtido de uma credencial validada, e não aceite como um valor arbitrário fornecido pelo cliente. Pedidos sem autenticação ou autorização válida não devem receber dados do backend.
    \item \textbf{Políticas explícitas por perfil:} cada política deve indicar os campos permitidos, mascarados, pseudonimizados, omitidos ou agregados, de acordo com a finalidade prevista. A configuração deve seguir o princípio do menor privilégio e evitar que campos novos sejam expostos implicitamente.
    \item \textbf{Proteção do fluxo e dos componentes:} devem ser consideradas a proteção das comunicações, a gestão de credenciais e segredos, a validação de entradas e respostas e a divulgação mínima de informação em erros e registos.
    \item \textbf{Verificação funcional e de segurança:} os testes devem confirmar que cada perfil recebe apenas a resposta prevista, que pedidos inválidos ou não autorizados são rejeitados e que as transformações não alteram indevidamente os dados necessários à finalidade do serviço.
    \item \textbf{Avaliação do impacto:} além de verificar as respostas por perfil, a avaliação prática deverá medir a latência, o consumo de CPU e memória e o tamanho do payload antes e depois da transformação. Estas métricas caracterizam o protótipo, mas não medem, isoladamente, o cumprimento das obrigações legais.
\end{itemize}

\noindent Deste modo, o RGPD orienta sobretudo a necessidade e a forma do tratamento de dados pessoais, enquanto a NIS2 enquadra a gestão de riscos de cibersegurança e a resiliência das entidades abrangidas. A arquitetura proposta pode contribuir para objetivos relacionados com ambos, mas a avaliação de conformidade exigiria analisar o tratamento completo, os responsáveis envolvidos, o contexto operacional e as medidas organizacionais e técnicas aplicáveis.

\chapter{Estado da Arte e Trabalhos Relacionados}

\section{Modelos Base}
\noindent A proposta deste projeto baseia-se na fusão e adaptação dos seguintes modelos de referência:
\begin{itemize}
    \item \text{Modelo PRIPRO / PRIADA}% Focado na gestão de perfis de privacidade do utilizador e na adaptação dinâmica da informação com base em permissões e contexto \cite{cesconetto2020}.
    \item \text{Modelo PADRES}% Propõe esquemas de redução de dados e encriptação consciente da privacidade (\textit{privacy-aware data reduction and encryption scheme}), promovendo a remoção ou agregação de atributos secundários para diminuir o tamanho dos dados transmitidos na rede \cite{pereira2022}.
\end{itemize}

\section{Análise dos Modelos}
\subsection{Modelo PRIPRO / PRIADA}

O modelo PRIPRO (\textit{PRIvacy PROfiles}) integra a arquitetura UbiPri (\textit{Ubiquitous Privacy}), desenvolvida para a gestão e controlo da privacidade em ambientes ubíquos. Esta arquitetura encontra-se organizada em módulos independentes, responsáveis por diferentes dimensões da gestão da privacidade \cite{leithardt2015,cesconetto2020,prichain2019}.

No contexto desta arquitetura, o módulo PRIPRO é responsável pela gestão dos perfis dos utilizadores, permitindo associar diferentes permissões e restrições aos recursos e serviços disponíveis. O objetivo é garantir que o acesso à informação e aos recursos do ambiente respeita as preferências e condições estabelecidas para cada perfil.

Uma das características relevantes do PRIPRO é a possibilidade de evolução dinâmica dos perfis. A abordagem considera a frequência com que um utilizador interage com determinado ambiente para gerir a progressão ou regressão do seu perfil. Esta evolução é suportada por uma taxonomia de métricas, parâmetros e critérios hierárquicos, permitindo adaptar as permissões às características da interação \cite{cesconetto2020}.

Complementarmente, o módulo PRIADA (\textit{PRIvacy ADAptation}) assume a responsabilidade pela gestão e controlo da adaptação, tratando informação relacionada com a adaptação de software e hardware no ambiente ubíquo. Enquanto o PRIPRO se concentra na gestão dos perfis e das respetivas permissões, o PRIADA está associado à adaptação do comportamento do sistema em função das condições identificadas. \cite{prichain2019,cesconetto2020,lopes2020priada}.

O PRIADA foca-se na gestão e adaptação da informação considerando a localização e o contexto da interação entre o indivíduo e o ambiente. O modelo procura determinar as condições de acesso à informação com base nas circunstâncias em que ocorre a interação, evidenciando a importância de considerar o contexto como uma variável na aplicação das políticas de privacidade \cite{lopes2020priada}.

\subsection{Modelo PADRES}
O PADRES (\textit{Tool for PrivAcy, Data REgulation and Security}), desenvolvido por Pereira, Crocker e Leithardt (2022), é uma ferramenta de código aberto concebida para apoiar a análise da privacidade e segurança de aplicações Web, tendo como principal enquadramento os requisitos do Regulamento Geral sobre a Proteção de Dados (RGPD). A sua motivação está relacionada com a complexidade do processo de avaliação da conformidade regulamentar, que exige a análise de procedimentos organizacionais, aplicações e infraestruturas tecnológicas \cite{pereira2022}.

A ferramenta organiza os principais princípios do RGPD sob a forma de listas de verificação e questionários, cujas respostas são introduzidas manualmente pelo utilizador. Esta abordagem permite sistematizar a avaliação de diferentes aspetos relacionados com o tratamento de dados pessoais, facilitando a identificação de possíveis lacunas relativamente aos requisitos regulamentares. Assim, o PADRES funciona como um instrumento de apoio à análise e à auditoria regulamentar \cite{pereira2022}.

Para complementar a componente regulamentar, o PADRES integra ferramentas externas de análise de segurança, nomeadamente o Nmap, o OWASP ZAP e analisadores de cookies. Os resultados obtidos através destas ferramentas são armazenados e posteriormente combinados com as respostas dos questionários, permitindo gerar um relatório consolidado com sugestões e recomendações. Esta integração procura reunir, num único processo, informação relativa à avaliação da privacidade e à identificação de potenciais vulnerabilidades de segurança nas aplicações analisadas \cite{pereira2022}.


\section{Análise Comparativa dos Trabalhos Relacionados}
\subsection{Trabalhos Relacionados}
Os trabalhos analisados abordam diferentes dimensões da proteção da privacidade, desde a gestão dos perfis dos utilizadores até à avaliação da conformidade regulamentar das aplicações. O PRIPRO e o PRIADA concentram-se na gestão e adaptação das políticas de privacidade em função dos perfis e do contexto, enquanto o PADRES procura apoiar a identificação de aspetos relacionados com a privacidade e a segurança de aplicações Web. Apesar de partilharem a preocupação com a proteção da informação, os seus objetivos e mecanismos de implementação são distintos \cite{cesconetto2020,lopes2020priada,pereira2022}.

Uma abordagem mais diretamente relacionada com a minimização de dados em interfaces de programação é apresentada por Pallas et al. (2022), através da solução Janus. Este trabalho aborda o problema da disponibilização de informação excessiva por APIs Web, propondo um mecanismo configurável de minimização de dados para APIs GraphQL. A solução permite definir diferentes operações de redução de informação de acordo com os papéis dos consumidores, procurando disponibilizar apenas os atributos necessários a cada pedido. Os autores apresentam ainda uma implementação experimental e uma avaliação do desempenho, demonstrando a viabilidade da aplicação de políticas de minimização diretamente nas interfaces de acesso aos dados \cite{pallas2022janus}.

Outra contribuição relevante é apresentada por Chen et al. (2022), através do sistema minTAP (\textit{Practical Data Access Minimization in Trigger-Action Platforms}). Esta abordagem procura reduzir o acesso excessivo a dados pessoais em plataformas que permitem estabelecer regras de automação entre diferentes serviços. Os autores identificam dois tipos de acesso excessivo: ao nível dos atributos, quando uma aplicação recebe mais campos do que aqueles de que necessita, e ao nível dos serviços, quando lhe são disponibilizadas mais APIs do que as necessárias. Para responder a este problema, o minTAP procura identificar os dados estritamente necessários à execução de cada regra, limitando a informação disponibilizada às aplicações consumidoras \cite{chen2022mintap}.

A análise conjunta destes trabalhos permite identificar duas dimensões complementares. Por um lado, os modelos PRIPRO e PRIADA demonstram a importância de utilizar perfis e informação contextual para determinar as condições de acesso e adaptação. Por outro, as abordagens Janus e minTAP demonstram a possibilidade de aplicar a minimização ao nível dos atributos disponibilizados às aplicações consumidoras. O PADRES acrescenta uma perspetiva de avaliação regulamentar e de segurança, permitindo enquadrar a necessidade de verificar os mecanismos de proteção implementados \cite{cesconetto2020,lopes2020priada,pereira2022,pallas2022janus,chen2022mintap}.

\subsection{Resumo da análise comparativa}

Apesar da existência destas abordagens, a proposta deste projeto procura reunir estas dimensões numa camada intermédia dedicada à adaptação de respostas HTTP/REST. Ao contrário de uma solução centrada exclusivamente na gestão de permissões, pretende-se que o middleware utilize o perfil associado ao pedido para determinar não só se o acesso é permitido, mas também que informação pode ser devolvida. A transformação poderá envolver a seleção de campos, a mascaragem ou pseudonimização de atributos e a agregação de informação, de acordo com as políticas definidas para cada perfil.

Deste modo, pretende-se explorar uma arquitetura que combine a gestão de perfis inspirada no PRIPRO/PRIADA com os princípios de minimização de dados presentes em trabalhos como Janus e minTAP, mantendo o enquadramento regulamentar e de segurança associado ao PADRES. A avaliação experimental permitirá analisar a eficácia funcional das transformações e o respetivo impacto na latência, no consumo de recursos e no volume de dados transmitidos. Esta combinação constitui a principal diferenciação da abordagem proposta relativamente aos trabalhos analisados.

\section{Tabela Comparativa dos Trabalhos Relacionados}
A Tabela \ref{tab} apresenta uma comparação dos trabalhos analisados na secção anterior, considerando os seus objetivos, mecanismos de proteção, utilização de perfis, adaptação da informação e aplicabilidade ao contexto das APIs. A comparação permite identificar as características já exploradas na literatura e as funcionalidades que se pretende reunir no middleware proposto.

Neste enquadramento, o middleware proposto procura integrar a gestão de perfis com a transformação direta das respostas de APIs REST, permitindo aplicar diferentes níveis de proteção aos mesmos recursos. A sua avaliação experimental permitirá verificar se esta integração consegue reduzir a quantidade de informação transmitida e a exposição de atributos sensíveis, mantendo a utilidade funcional das respostas e introduzindo um custo computacional aceitável.

%%tabela%%
\begin{table}[H]
\centering
\caption{Comparação dos modelos e trabalhos relacionados}
\label{tab}
\renewcommand{\arraystretch}{1.4}
\setlength{\tabcolsep}{3pt}
\scriptsize

\resizebox{\textwidth}{!}{%
\begin{tabular}{|p{2.2cm}|p{2cm}|p{2.3cm}|p{1.6cm}|p{2.5cm}|p{2.0cm}|p{2.0cm}|}
\hline
\textbf{Modelo / Abordagem}
& \centering\textbf{Controlo de Acesso}
& \centering\textbf{Mascaragem}
& \centering\textbf{Redução ($\Delta D$)}
& \centering\textbf{Adaptabilidade}
& \centering\textbf{Minimização de Dados}
& \centering\textbf{Âmbito de Aplicação}
\tabularnewline
\hline

PRIPRO / PRIADA
& Perfis e contexto
& Não especificado
& Não especificado
& Dinâmica
& Não especificada
& Ambientes ubíquos
\tabularnewline
\hline

PADRES
& Não aplicável
& Não
& Não
& Avaliação e recomendações
& Não implementada
& Aplicações Web
\tabularnewline
\hline

Janus
& Baseado em papéis
& Não especificado
& Sim
& Configurável
& Seleção de atributos
& APIs GraphQL
\tabularnewline
\hline

minTAP
& Dados e serviços
& Não especificado
& Sim
& Dependente das regras
& Atributos e serviços
& Plataformas Trigger-Action
\tabularnewline
\hline

\textbf{Proposta (Middleware)}
& \textbf{Contextual / Perfis}
& \textbf{Sim}
& \textbf{Sim}
& \textbf{Dinâmica}
& \textbf{Seleção, mascaragem e agregação}
& \textbf{APIs: HTTP e REST}
\tabularnewline
\hline
\end{tabular}}
\end{table}
%%

\chapter{Arquitetura Proposta para o Middleware}
\noindent O sistema atua como um proxy de segurança em Python intercetando pedidos entre clientes e serviços de backend.

\section{Perfis de Acesso Adaptativos}
\noindent O middleware processa as respostas dinamicamente segundo três perfis principais:
\begin{itemize}[leftmargin=*]
    \item \textbf{Perfil Administrador:} Recebe o conjunto de dados completo com permissões integrais de administração.
    \item \textbf{Perfil Operacional / Analista:} Recebe os dados com mascaragem automática e pseudonimização de atributos confidenciais.
    \item \textbf{Perfil Público / Terceiros:} Recebe apenas dados agregados e métricas consolidadas, minimizando o payload enviado.
\end{itemize}

\chapter{Conclusões e Trabalhos Futuros}
\noindent A formalização teórica apresentada neste relatório valida a viabilidade de combinar o controlo adaptativo por perfis (PRIPRO) com o esquema de redução de dados consciente da privacidade (PADRES). Como trabalhos futuros no âmbito da Parte 2 do projeto, proceder-se-á à implementação do protótipo em Python recorrendo às frameworks FastAPI/Flask, desenvolvendo scripts de benchmarking para medir latência (ms), consumo de CPU e memória.

\renewcommand{\bibname}{Referências Bibliográficas}
\bibliographystyle{unsrt}
\bibliography{references}

\end{document}